\appendix
\section{Analytical Justification of Regime Thresholds}
\label{app:regime}

This appendix provides an analytical justification for the regime thresholds
$R_C \approx \theta_C$ and $R_B \approx \theta_B$ introduced in the main text.
Our goal is not to derive exact constants, which are inherently platform- and
implementation-dependent, but to show that the existence of critical inflection
points is structurally inevitable in hierarchical memory systems.

\subsection{Latency lower bounds in hierarchical memory systems}
We begin from the latency decomposition used in the main text:
\begin{equation}
T_{\text{total}} = T_{\text{comp}} + T_{\text{mem}} + T_{\text{swap}} + T_{\text{sync}}.
\end{equation}
When the active working set of size $W$ exceeds fast-memory capacity $C_{\text{fast}}$,
locality-driven access costs impose a lower bound
$T_{\text{mem}} \gtrsim \alpha \, W / C_{\text{fast}}$, where $\alpha > 0$ captures
platform-dependent miss penalties. Independently, any compulsory cross-tier transfer
of data volume $D$ across the slowest tier with sustained bandwidth $B_{\text{slow}}$
incurs an irreducible cost $T_{\text{swap}} \gtrsim \beta \, D / B_{\text{slow}}$.

These bounds hold regardless of scheduling or overlap: while overlap can hide portions
of transfer latency, it cannot violate physical bandwidth or residency constraints.

\subsection{Dimensionless ratios and inflection points}
Introducing the dimensionless ratios
\begin{equation}
R_C = \frac{C_{\text{fast}}}{W}, \quad
R_B = \frac{B_{\text{slow}}}{D},
\end{equation}
we observe that when either $R_C$ or $R_B$ falls below an $\mathcal{O}(1)$ constant,
the corresponding lower-bound term dominates $T_{\text{total}}$.
At these points, coordination and synchronization overheads transition from being
amortizable to additive, producing an abrupt change in the dominant latency component.

This inflection behavior explains why regime transitions are discontinuous rather than
gradual. Small changes in workload size, batch configuration, or sustained bandwidth
can move the operating point across a threshold where locality or transfer constraints
become dominant.

\subsection{Interpretation of threshold values}
The specific threshold values $\theta_C \approx 0.5$ and $\theta_B \approx 0.4$ reported
in the main text reflect empirical operating points at which the lower-bound terms
begin to dominate in our experimental platforms. While these numerical values are not
universal constants, their order and existence are robust: any hierarchical memory
system must exhibit similar inflection points when capacity or bandwidth constraints
are approached.

Thus, the regime thresholds should be interpreted as structural boundaries induced by
physical constraints, not as tunable parameters. This analytical perspective supports
the central claim of this work that regime-dependent behavior is an intrinsic property
of hierarchical memory orchestration.